% =====================================================================
%  Style dùng chung cho mọi hình TikZ/pgfplots của báo cáo.
%  Bảng màu phong phú, khối có bóng đổ và chuyển sắc, macro khối hộp 3D.
% =====================================================================
\usepackage{tikz}
\usepackage{pgfplots}
\pgfplotsset{compat=1.18}
\usetikzlibrary{external,positioning,arrows.meta,matrix,fit,calc,shapes.geometric,shapes.misc,shapes.arrows,
                backgrounds,decorations.pathreplacing,decorations.pathmorphing,shadows,shadows.blur,patterns,3d,
                perspective,fadings,shadings,trees}
\usepgfplotslibrary{groupplots,colormaps,fillbetween}

% Externalize: mỗi hình biên dịch một lần thành tikzcache/<tên>.pdf.
\tikzexternalize[prefix=tikzcache/]
\newcommand{\hinh}[1]{\tikzsetnextfilename{#1}\input{tikz/#1}}

% Cỡ chữ trong hình: lớp 14pt cho \tiny = 6pt, \scriptsize = 8pt, \footnotesize = 10pt; hình lại bị co theo bề
% rộng trang nên chữ còn nhỏ hơn. Trong mọi tikzpicture nâng ba cỡ này lên để chữ hình gần cỡ chữ thân bài.
\makeatletter
\tikzset{every picture/.append style={execute at begin picture={%
  \def\tiny{\@setfontsize\tiny{9.5}{11.5}}%
  \def\scriptsize{\@setfontsize\scriptsize{11}{13}}%
  \def\footnotesize{\@setfontsize\footnotesize{12}{14}}}}}
\makeatother

% ----------------------------------------------------------------- bảng màu
\definecolor{mZero}{HTML}{7F7F7F}
\definecolor{mOne}{HTML}{0072B2}
\definecolor{mTwo}{HTML}{E69F00}
\definecolor{mThree}{HTML}{009E73}
\definecolor{accentRed}{HTML}{C0392B}
\definecolor{lightBlue}{HTML}{DCE8F5}
\definecolor{lightOrange}{HTML}{FCE9CF}
\definecolor{lightGreen}{HTML}{D5F0E6}
\definecolor{lightGray}{HTML}{EDEDED}
% màu tươi cho sơ đồ
\definecolor{cBlue}{HTML}{2F80ED}
\definecolor{cSky}{HTML}{56CCF2}
\definecolor{cTeal}{HTML}{1ABC9C}
\definecolor{cGreen}{HTML}{27AE60}
\definecolor{cLime}{HTML}{A3D977}
\definecolor{cYellow}{HTML}{F2C94C}
\definecolor{cOrange}{HTML}{F2994A}
\definecolor{cRed}{HTML}{EB5757}
\definecolor{cPink}{HTML}{E86FA8}
\definecolor{cPurple}{HTML}{9B51E0}
\definecolor{cIndigo}{HTML}{5B5FE0}
\definecolor{cBrown}{HTML}{A0522D}
\definecolor{cNavy}{HTML}{123A75}
\definecolor{cSlate}{HTML}{4F5D75}

\pgfplotsset{
  a06/.style={
    width=8cm, height=5.6cm, grid=major, grid style={gray!25},
    tick label style={font=\small}, label style={font=\normalsize},
    legend style={font=\small, draw=gray!50, fill=white, fill opacity=0.9, text opacity=1, rounded corners=2pt},
    title style={font=\normalsize\bfseries},
    /pgf/number format/use comma, /pgf/number format/1000 sep={.}},
  m0/.style={mZero, very thick, mark=none},
  m1/.style={mOne, very thick, mark=none},
  m2/.style={mTwo, very thick, mark=none},
  m3/.style={mThree, very thick, mark=none},
}

\tikzset{
  % khối chữ nhật bo góc có bóng mờ và chuyển sắc từ trên xuống
  hop/.style={rounded corners=3pt, draw=#1!75!black, thick, top color=#1!12, bottom color=#1!38,
              blur shadow={shadow blur steps=6, shadow xshift=0.8pt, shadow yshift=-0.8pt, shadow opacity=25},
              align=center, inner sep=4pt, font=\small},
  hop/.default=cBlue,
  khoi/.style={hop=cBlue, minimum height=8mm},
  khoiC/.style={hop=cOrange, minimum height=8mm},
  khoiG/.style={hop=cGreen, minimum height=8mm},
  khoiX/.style={hop=cSlate, minimum height=8mm},
  khoiP/.style={hop=cPurple, minimum height=8mm},
  khoiR/.style={hop=cRed, minimum height=8mm},
  khoiT/.style={hop=cTeal, minimum height=8mm},
  % quả cầu có bóng (nơ-ron, phép toán)
  cau/.style={circle, inner color=white, outer color=#1!55, draw=#1!65!black, thick, text=black!85,
              font=\small\bfseries, inner sep=1pt, minimum size=8mm,
              blur shadow={shadow blur steps=5, shadow opacity=30}},
  cau/.default=cBlue,
  mui/.style={-{Stealth[length=2.6mm]}, very thick, cNavy},
  muiM/.style={-{Stealth[length=2.6mm]}, very thick, #1},
  nhan/.style={font=\small, text=black!75},
  o/.style={draw=gray!70, minimum size=7mm, inner sep=0pt, font=\small},
  tieude/.style={font=\normalsize\bfseries\sffamily, text=cNavy},
}

% \khoibaD{x}{y}{rộng}{cao}{sâu}{màu}: khối hộp 3D, góc dưới-trái-trước tại (x,y),
% chiều sâu vẽ xiên lên phải (0,6 × sâu theo 30 độ). Mặt trên sáng, mặt bên tối.
\newcommand{\khoibaD}[6]{%
  \pgfmathsetmacro{\kdx}{0.6*#5*cos(30)}\pgfmathsetmacro{\kdy}{0.6*#5*sin(30)}%
  \filldraw[draw=#6!60!black, fill=#6!85!black, line join=round] (#1,#2) rectangle ++(#3,#4);
  \filldraw[draw=#6!60!black, fill=#6!45, line join=round] (#1,#2+#4) -- ++(\kdx,\kdy) -- ++(#3,0) -- ++(-\kdx,-\kdy) -- cycle;
  \filldraw[draw=#6!60!black, fill=#6!60!black, line join=round] (#1+#3,#2) -- ++(\kdx,\kdy) -- ++(0,#4) -- ++(-\kdx,-\kdy) -- cycle;
}
% \tensorbaD{x}{y}{rộng}{cao}{sâu}{màu}{nhãn}: khối 3D có nhãn phía dưới (feature map, tensor, lớp mạng).
\newcommand{\tensorbaD}[7]{%
  \khoibaD{#1}{#2}{#3}{#4}{#5}{#6}%
  \node[font=\small, text=black!80, align=center, anchor=north] at ({#1+0.5*#3},{#2-0.08}) {#7};
}

\tikzset{
  % trụ cơ sở dữ liệu / kho tri thức
  csdl/.style={cylinder, shape border rotate=90, aspect=0.25, draw=#1!70!black, thick,
               top color=#1!15, bottom color=#1!45, minimum width=16mm, minimum height=14mm, align=center,
               font=\small, blur shadow={shadow blur steps=5, shadow opacity=25}},
  csdl/.default=cTeal,
  % nhãn tròn nhỏ đánh số bước
  buoc/.style={circle, fill=#1, text=white, font=\small\bfseries\sffamily, inner sep=1.5pt, minimum size=6mm},
  buoc/.default=cRed,
  % vùng nền nhóm khối
  vung/.style={rounded corners=6pt, draw=#1!60, dashed, thick, fill=#1!6, inner sep=8pt},
  vung/.default=cSlate,
}
